\section{Intro}

Now for the new stuff.
We first bring the Focked up ZX-calculus up to speed with the qubit ZX-calculus, essentially.
We start by defining Pauli webs, and determine the general form of a Pauli web on a diagram in the Gaussian+combs fragment.
Then we define what it means for two such diagrams to be fault-equivalent, and we give a slew of fault-equivalent rewrite rules.
These are put to work as we give two applications, both of which are analogues of things we've discussed in the qubit case already.

We first show that the ZX-calculus can move us between static and dynamic implementations of GKP codes, just like in \Cref{ch:floquetifying}, which itself was based on \Ccite{Bombin_2024}.
Following the latter, we also include a \textit{measurement-based} implementation in this qumode case.
Most importantly, unlike in \Cref{ch:floquetifying}, in this chapter we use fault-equivalent rewrite rules to do this, so we know we aren't decreasing the performance of the code under the ZX noise model through this process.
To our knowledge, this is the first discussion of dynamic codes in the CV setting.

Next, we use fault-equivalent rewrite rules to come up with a syndrome extraction circuit for a particular measurement in a multi-mode GKP code.
Again, to our knowledge, there's very little discussion of syndrome extraction circuits for multimode GKP codes in the presence of realistic noise models in the literature, so we hope we are adding to an underdeveloped area here.
